\chapter*{Preface}

In 1936 a young mathematician answered a question of Hilbert by imagining a clerk with a pencil, an eraser and an
endless strip of paper, and proved that some questions can never be settled by any such clerk, however patient. In
1931 another had shown that some true statements about whole numbers can never be proved. Both arguments are
descendants of Cantor's diagonal, and both turned a vague hope (that every precise question has a mechanical answer)
into a theorem about its failure. This book tells what came of these two discoveries: a theory of what can be
computed and proved at all, and then, once computers existed, a theory of what can be computed \emph{quickly}.

We begin with the simplest machines, finite automata, which already have a complete and beautiful theory, and climb
through grammars and pushdown automata to Turing machines. Then we meet the undecidable: the halting problem,
Post's correspondence problem, tilings of the plane, and Hilbert's tenth problem on Diophantine equations. Two
chapters turn to logic: the $\lambda$-calculus, where programs and proofs turn out to be the same objects, and
G\"odel's theorems, completeness first and incompleteness second, with Goodstein's theorem as an example of a true
statement about numbers that ordinary arithmetic cannot prove. The second half is complexity theory: time and space
and their hierarchies, NP-completeness and the reductions that show thousands of problems to be equally hard, the
polynomial hierarchy and counting, the power of randomness and of interaction, and finally circuits and proofs, the
places where lower bounds have actually been proved, and the barriers that explain why so few have.

The prerequisites are Tomes~C1 and~C2, with a little of Tomes~1 and~3; results from them are recalled where they are
used, with a precise reference. Every theorem is proved, except for a few landmarks whose proofs would need a book
of their own; those are clearly labelled as quoted, with sources. Each chapter opens with a question to think about,
and ``Pause and think'' boxes ask you to predict before reading on. Exercises are graded from $\star$ to
$\star\star\star$; each has a hint and a solution that begins by explaining how one might find the idea.

Each chapter has a companion Python notebook that builds the objects of the chapter and runs them: automata and
their minimisation, Turing machines and busy beavers, $\lambda$-terms reducing themselves, G\"odel numbers, the
reductions between NP-complete problems checked by brute force, interactive protocols, and the smallest circuits
for small functions. The negative theorems of the book become experiments: a reader can watch a search for a
solution of Post's problem run on and on, and understand why nothing can tell in advance whether it will stop.
